% LSPV-3 REPORT::ORBITDECODERFW1::2026-10-13T23:36:00+09:00
% !TEX program = xelatex
\documentclass[11pt]{article}
\usepackage[T1]{fontenc}
\usepackage{lmodern,amsmath,amssymb,amsthm,booktabs,tabularx,xurl,hyperref}
\usepackage[margin=0.84in]{geometry}
\newtheorem{theorem}{Theorem}
\newtheorem{corollary}[theorem]{Corollary}
\newcommand{\F}{\mathbb F_2}
\newcommand{\EXACT}{\textbf{PROVED/EXACT}}
\newcommand{\KNOWN}{\textbf{KNOWN THEOREM}}
\newcommand{\COMP}{\textbf{COMPUTATIONAL}}
\newcommand{\COND}{\textbf{CONDITIONAL}}
\newcommand{\CAND}{\textbf{CANDIDATE}}
\newcommand{\OPEN}{\textbf{OPEN}}
\newcommand{\RET}{\textbf{FALSIFIED/RETIRED}}
\hypersetup{colorlinks=true,urlcolor=blue,linkcolor=blue}
\setlength{\emergencystretch}{3em}
\title{ORBITDECODERFW1: Automorphism Obstruction to\\Untagged Source Orientation Is Not Hardness}
\author{UNIVERSAL\_CORE\_REDTEAM / LSPV-3}
\date{13 October 2026, 23:36 KST}

\begin{document}
\maketitle

\begin{center}
\fbox{\parbox{0.95\linewidth}{
\textbf{Status: P versus NP remains OPEN.}
The explicit decoded hierarchy theorem survives.  This run isolates its exact
scope boundary.  An untagged unique-extension relation can have an
automorphism exchanging the designated source and extension blocks, so no
variable-renaming-covariant target-only decoder can recover the original
orientation on every instance.  The smallest nonlinear control is the graph
of the quadratic involution
$S(a,b,c)=(a,b,c\oplus ab)$.  Yet $S^{-1}=S$, so the same example is uniformly
easy in both directions.  Hence hiding source labels is \RET{} as a
sufficient hardness repair.  The orbit theorem has a complete proof but is
\CAND{} pending independent review; the bounded audit is \COMP{}.}}
\end{center}

\section{Claim card and exact quantifiers}
The claim attacked is the implicit extension of HIERLIFT-TEST1 from its
declared record encoding to arbitrary untagged formula presentations:
\begin{equation}
\begin{aligned}
&\exists\text{ uniform polynomial }D\ \forall(C,X)\ \forall\pi:\\
&\qquad D(\pi C)=\pi D(C),\qquad D(C)=X,
\end{aligned}
\tag{1}
\end{equation}
where $C$ is a unique-extension relation, $X$ is its originally designated
source-variable block, and $\pi$ is a variable renaming.  Clauses are treated
as an unordered set.  Equation (1) asks $D$ to be target-only and
renaming-covariant; it receives no sidecar designation of $X$.

The exact counterexample prefix is
\[
\exists(C,X,\tau)\ \forall D\;[
 \tau C=C\land \tau X\ne X\land
 D(\tau C)=\tau D(C)\Rightarrow D(C)\ne X].              \tag{2}
\]
The machine is a deterministic uniform random-access decoder with arbitrary
polynomial workspace.  The SAME\_REALIZATION tuple is the same canonical CNF
$C$, its same automorphism $\tau$, and the same requested source block $X$.
The weakest material result is a counterexample to (1), not a runtime lower
bound.

This distinction separates two uses of a reduction.  An external
source-to-target wrapper may retain the source and its construction metadata.
In contrast, a short program counted by $H(G)$ receives the target $G$; if it
plans to run a source program, it must recover the source from $G$ or from a
declared tag within the same charged input.

\section{Primary-literature and cross-domain preflight}
Crawford, Ginsberg, Luks, and Roy study adding symmetry-breaking predicates to
search problems \cite{CGLR}.  Their work is a useful structural analogy:
symmetry can remove a canonical orientation and explicit predicates can
restore one.  It does not imply a lower bound for arbitrary witness search.
Babai's quasipolynomial graph/string-isomorphism algorithm \cite{Babai}
reinforces the algorithmic warning that canonization and group action are
their own computational tasks; absence of a positional parser is not evidence
of SAT hardness.

Berman and Hartmanis study polynomial-time isomorphisms of complete sets
\cite{BH}.  The directly transferable part here is only elementary: an
explicit signed permutation with a known inverse is a polynomial-time
isomorphic recoding, so a wrapper that retains it transports witnesses at
polynomial cost.  Their general NP-isomorphism conjectural architecture is
not imported.  Thus the literature classifications are:
\begin{itemize}
\item explicit invertible recoding: \textbf{THEOREM TRANSFER} after the map
and its inverse are actually supplied;
\item symmetry-breaking/canonization as a repair: \textbf{ANALOGY/CANDIDATE}
until a decoder for the declared target encoding is proved;
\item symmetry or canonization difficulty $\Rightarrow$ unrestricted SAT
time: \textbf{FAILED TRANSFER}.
\end{itemize}

\section{Orbit obstruction theorem}
Let a CNF automorphism be a variable permutation preserving its canonical
unordered clause set.  A decoder $D$ is covariant when
$D(\pi C)=\pi D(C)$.

\begin{theorem}[automorphism-orbit obstruction]
If $\tau C=C$ but $\tau X\ne X$, then no covariant decoder can output the
designated block $X$ on $C$.
\end{theorem}
\begin{proof}
Covariance and $\tau C=C$ give
\[
 D(C)=D(\tau C)=\tau D(C).                               \tag{3}
\]
If $D(C)=X$, equation (3) says $X=\tau X$, contrary to the hypothesis.
Therefore a covariant target-only decoder cannot recover the non-invariant
designation.
\end{proof}

The proof is complete and elementary, but this newly recorded project claim
is retained as \CAND{} until independent review under the promotion rule.

\section{Nonlinear unique-extension counterexample}
Define the quadratic shear
\[
 S:\F^3\to\F^3,\qquad S(a,b,c)=(a,b,c\oplus ab).          \tag{4}
\]
It is an involution: $S(S(x))=x$.  Let
\[
 R(x,y)\iff y=S(x).                                      \tag{5}
\]
Encode (5) canonically by one width-six clause forbidding each invalid
assignment.  There are $64-8=56$ distinct proper clauses and eight models.

\begin{corollary}[nonlinear orientation failure]
The block swap $\tau(x,y)=(y,x)$ is an automorphism of $R$, but sends the
designated source block $X$ to the disjoint target block $Y$.  Thus Theorem 1
rules out a covariant target-only decoder that always recovers $X$.
Nevertheless every $x$ has the unique extension $S(x)$, every $y$ has the
unique predecessor $S(y)$, and both maps take constant circuit depth and
linear output work.
\end{corollary}
\begin{proof}
By involutivity,
\[
 y=S(x)\iff x=S(y),                                      \tag{6}
\]
so swapping blocks preserves exactly the model relation and hence its
canonical forbidden-assignment CNF.  The unique maps in both directions are
$S$ itself.  The formula contains one AND and one XOR beyond copied bits.
\end{proof}

The corollary defeats two opposite overstatements.  Unique extension does not
provide a canonical source orientation, but failure of orientation does not
provide unavoidable work.  Equality paths give an affine Krom version: path
reversal exchanges the endpoints even though either endpoint uniquely fixes
the whole path.

\section{Exact repair and effect on HIERLIFT/PERMEXT}
HIERLIFT-TEST1 uses a record beginning with the literal header
\texttt{HL2\textbackslash x00}, source width, depth, phase, field records, and
acyclic dependencies.  Its decoder theorem is unaffected.  PERMEXT1 likewise
declares source variables, graph, cotree, field polynomial, circuits, and
topological gate order.  At that structured-encoding scope, the unique Gold
extension and its polynomial composition remain \EXACT{} after the preceding
independent review.

For a standalone untagged CNF, one of the following must be stated:
\begin{enumerate}
\item a tag/sidecar is part of the legal encoding and invalid records are
rejected in polynomial time;
\item a polynomial canonical decoder is proved for that particular image;
\item the algorithm is an external reduction wrapper retaining the source and
all permutation/phase metadata; or
\item no source orientation is required because every orbit representative
admits a uniform extension algorithm.
\end{enumerate}
Options 1, 3, and 4 are algorithmic escapes, not lower bounds.  Option 2 is a
new theorem obligation.  A forward many-one reduction alone neither grants
nor forbids a target-only inverse decoder.

\section{Bounded reproducible audit}
The dependency-free checker
\path{experiments/review/decoder_orbit_symmetry_firewall.py} used seed
\texttt{23362026}; its SHA-256 is
\texttt{3be2d56eca1079e484caac4761750c7b355ee6d0cda686508997d501473bc106}.
One invocation completed in 1.128096 seconds, so no chunking was needed.

It passed 1,156 assertions.  The audit exhausted all 64 assignments of the
six-variable shear relation and all $6!=720$ variable permutations.  It found
eight models, a 16-element CNF automorphism group, and an eight-element orbit
of the designated three-variable source block.  It also enumerated every
automorphism of equality paths of lengths two through eight; each path had the
two expected automorphisms and the source endpoint orbit contained both
endpoints.

Across 512 seeded signed-variable renamings with independent clause shuffles,
the metadata-aware inverse recovered all 512 original witnesses.  A naive
first-three-variables parser happened to preserve the true transported source
block only 35 times.  Six malformed metadata records were rejected.  YES,
NO, Horn, Krom, affine, bounded-path, independent-product, and planted controls
were evaluated.  The immutable raw record is
\path{results/raw/review/2026-10-13_2336_decoder_orbit_symmetry_firewall.json},
SHA-256
\texttt{59ea6b1fe17c7b353cb3a51c1c14f54d26a6c4fa61f52b4dd2e22d4cdeda1790}.
Finite checks are \COMP{} and are not extrapolated to asymptotic time.

\section{Hostile controls, barriers, and post-quantum ledger}
The shear relation is nonlinear, bijective, source-matched, and compact; the
equality paths are affine Krom and bounded treewidth.  The audit includes
renaming, complementation, clause shuffling, malformed inputs, YES/NO, Horn,
Krom, affine, path, product, and planted controls.  DEAD, COPY, DELAY,
INTERP, and RECOMP wrappers cannot turn the elementary orbit identity into a
lower bound; they either preserve the symmetry or add explicit decodable
structure.  A semantic oracle can output a witness directly, and an algebraic
oracle evaluates the shear, so the result relativizes and algebrizes.  No
Natural-Proofs property is used.

For the post-quantum ledger, the graph of a quadratic involution and its
automorphism orbit are an \textbf{EXACT OBJECT/THEOREM TRANSFER} from
quadratic-map structure.  Collision entropy is maximal-lossless because every
fiber is a singleton; it still does not select an orientation.  Treating two
openings related by the involution as a forbidden collision is a
\textbf{FAILED TRANSFER}: both openings are valid and efficiently related.
Schur squares, relation-rank spectra, Hilbert functions, and syzygies are
\textbf{ANALOGY/CANDIDATE} diagnostics for different residuals and supply no
decoder or unrestricted lower bound here.

\section{CORPUS REVIEWED}
The current GitHub tree was surveyed before substantive work.  This run read
root \path{main.tex}, \path{AGENTS.md}, \path{CODEX.md}, the complete
pnp-research skill and red-team/protocol references, all seven task
specifications, current \path{main4.tex} and \path{main5.tex}, the verified
AFFINE-RESIDUAL-REDIRECT1 directive, EARLY22/SINGLE22 scope guards, newest MAIN
CUBEPERMFW1, and the actual RANKNONLINFW1, COTREELIFT1, BIJECTIVEQFW1,
QUADSHEAR1, POLARKERNEL1, FIELDTRACELIFT1, PERMEXT1, and HIERLIFT-TEST1
reports and cited code/results.

The HEC/VSC-HEC/G-PPC-HEC/CRS-PPC/CCSPC branch was checked.  Its same warning
persists: hiding a schedule, pointer, source orientation, or decomposition
does not constrain an arbitrary static solver without a model-preserving
compiler.  Existing exact-$(2,2)$ compilers and Resolution/PC/PCR projections
remain preserved only at their declared scope.

The immediate predecessor
\path{research/agents/lspv-2/2026-10-13_2324_LSPV-2_HIERLIFT-TEST1.tex}
was read back from GitHub at blob
\texttt{649adcce18f51ecaf63339c0112ddb76ab913fd9}; it is substantive and has
the required header and sentinel.  No predecessor recovery request was due.

\section{DUPLICATION / CONTRADICTION AUDIT}
Repository search found earlier covariance warnings in SYNTHESIS-0027,
explicit decoder requirements in HIERLIFT-TEST1, p-isomorphic wrapper
controls, and the shear inversion theorem.  The new component is their exact
intersection: a nonlinear involutive unique-extension relation whose own
automorphism makes original source orientation undefinable by a covariant
target-only decoder.

There is no contradiction with HIERLIFT-TEST1 because its header and layer
order are part of the input.  There is no contradiction with PERMEXT1 because
its declared structured record exposes the needed fields.  There is no
contradiction with NP-completeness: the counterexample is deliberately easy.
The result does not generalize to every untagged Gold/Tseitin CNF; its exact
scope is the proposed universal decoder extension in (1).

\section{WHAT CHANGED RELATIVE TO ROOT MAIN}
Root main correctly promotes the explicit decoded hierarchy and says that an
arbitrarily renamed untagged CNF needs a decoder theorem.  This run sharpens
``needs a theorem'' to a formal impossibility for any universal covariant
decoder: automorphisms can exchange source and extension blocks.  It also
closes the tempting hardness repair in the opposite direction: the same
quadratic shear is uniformly invertible, so decoder ambiguity alone supplies
no work factor.  The whole-machine implication remains \OPEN{}.

\section{ADOPTED / REJECTED / RETIRED PRIOR COMPONENTS}
\begin{itemize}
\item \textbf{Adopted:} the explicit tagged HIERLIFT decoder, PERMEXT unique
extension, cotree witness bijection, and charged signed-permutation metadata.
\item \textbf{Preserved:} all EARLY22/SINGLE22 compiler and restricted proof
results within their original hypotheses.
\item \textbf{Rejected:} free recovery of an original source block from an
untagged automorphic target, hidden nonuniform orientation advice, and
canonization-to-runtime inference.
\item \textbf{Retired:} ``hide the source/layer labels'' as a sufficient
repair for residual hardness.  Either metadata transports the easy lift, or
orientation ambiguity still does not prove witness-search cost.
\end{itemize}

\section{STALE OR OVERSTATED CLAIMS FOUND}
It would be stale to describe PERMEXT1 as review-pending at its declared
structured scope after HIERLIFT-TEST1's independent audit.  It would be
overstated to extend that exact result automatically to arbitrary untagged
CNF presentations.  Conversely, it would be overstated to require a
target-only decoder for an external reduction that retains the source and
charged construction metadata.  These are different machine inputs and must
not be conflated.

\section{Status ledger and next target}
\begin{center}\small
\begin{tabularx}{\textwidth}{@{}l l X@{}}
\toprule
Component & Status & Scope\\
\midrule
Explicit decoded hierarchy & \EXACT & Declared tagged record and acyclic order.\\
PERMEXT structured Gold lift & \EXACT & Preceding independent review; polynomial wrapper.\\
Orbit-obstruction theorem & \CAND & Complete proof; independent review pending.\\
1,156-assertion audit & \COMP & Finite verification only.\\
Universal untagged covariant orientation & \RET & Block-swap automorphism counterexample.\\
Decoder ambiguity as hardness & \RET & Same shear has explicit inverse.\\
All-program nonlinear-residual lower bound & \OPEN & No model-matched implication.\\
P versus NP separation & \OPEN & Load-bearing arrow remains open.\\
\bottomrule
\end{tabularx}
\end{center}

\textbf{Next target.}  Stop using hidden orientation as the obstruction.
Specify a tagged, polynomially recognizable, source-matched residual whose
satisfiable image is guaranteed but for which no uniform source-witness
extension is known.  Attack it first with both orientations, automorphism
orbits, field inverses, shear/Feistel decompositions, affine equivalence,
quotient/fiber algorithms, tensor/direct-sum decompositions, and charged
Gr\"obner/linearization methods.  Only a model-matched theorem excluding all
eligible programs would advance the UWI lower bound.

\begin{thebibliography}{9}
\bibitem{CGLR}
J. Crawford, M. Ginsberg, E. Luks, and A. Roy,
``Symmetry-Breaking Predicates for Search Problems,''
\emph{Proceedings of KR'96}, 1996, pp.~148--159,
\url{https://kr.org/proceedings/KR-1996-proceedings-scanned.pdf}.

\bibitem{Babai}
L. Babai, ``Graph Isomorphism in Quasipolynomial Time,''
\emph{Proceedings of STOC 2016}, pp.~684--697,
\url{https://arxiv.org/abs/1512.03547}.

\bibitem{BH}
L. Berman and J. Hartmanis,
``On Isomorphisms and Density of NP and Other Complete Sets,''
\emph{SIAM Journal on Computing} 6(2), 1977, pp.~305--322,
\url{https://doi.org/10.1137/0206023}.
\end{thebibliography}

\end{document}
% END_OF_VERIFIED_REPORT::LSPV-3::ORBITDECODERFW1
